\chapter{Symbolic Execution and Rewriting for Vector Instructions}
\label{ch:symex}

When an instruction is symbolically executed, we need a way of interpreting what computation is actually being done. There are 2 principled ways to do this for any (nontrivial) instruction, each of which abstracts details to one place or another. The first way is to break the instruction down into a composition of simpler operations as soon as it is symbolically executed, and the second is to include a single atomic operation in the execution graph.

% approach 1
The first approach breaks vector instructions into their component operations in symbolic execution. Most vector instructions follow the same pattern: slice lane from both source registers, perform a binary operation, write to a lane of the destination register. For example, the \verb|vpaddq| (vector packed addition) instruction performs a 64-bit addition on each lane of two 256-bit YMM registers and writes the results to the corresponding lanes of the destination register. This entire instruction can be expressed as a composition of simpler operations: slicing, scalar addition, and \verb|set_slice| to combine results back into a vector.

This pattern can be captured by a few helper functions. The function doing the heavy-lifting takes a scalar binary operation and applies it to each lane of the input vectors, building up the result with \verb|set_slice| operations:

\codeshot{vector_binop_aux.png}

This makes the symbolic execution of most instructions simple to implement. For example, \verb|vpaddq| can be implemented as follows:

\codeshot{SymexNormalInstruction.png}

But this approach makes the symbolic execution graph much more complex, because a 4-lane vector instruction uses roughly 12 nodes to slice inputs, perform lane ops, and slice the output together. This has two downsides:
First, canonicalizing the execution graph becomes more complex. We can easily end up with towers of nested slices and \verb|set_slice| deep in the graph, and we need powerful rewrite rules to recognize patterns and read many layers into existing nodes to find them.
Second, the execution graph has lost information that might be useful for other reasons. 
\todo{expand}
% something here about synthesis, cryptopt? having a vadd node can hint to some other process to use vector addition... but i dont think this dag actually gets used in synthesis?


% approach 2
The second approach keeps symbolic execution simple, emitting a single atomic \verb|op| for each vector instruction and abstracting details away. 

This means we need to define new symbolic operations:
%
% Variant op := ... | vadd (lane_width num_lanes : N) | vsub (lane_width num_lanes : N).
%
and symbolic execution is now very simple:
%   | vpaddq, [dst; src1; src2] =>
%     let num_lanes := (s / 64)%N in
%     v1 <- GetOperand src1;
%     v2 <- GetOperand src2;
%     result <- App ((vadd 64 num_lanes), [v1; v2]);

But we now need to handle vadd op in \verb|interp_op|, which defines how symbolic ops correspond to actual machine state operations. like with the first approach, this is easily modularized because of the slice, binop, \verb|set_slice| pattern which vector operations all follow.

We define a helper function similar to the one used in the first approach:
%
% (* Lane-parallel vector interpretation: applies a scalar binary operation
%    independently to each lane of two packed vectors, combining results. *)
% Fixpoint interp_vector_binop (scalar_op : Z -> Z -> Z) (lane_width : Z)
%   (lane_idx num_remaining : nat) (a b : Z) : Z := ...
%
% to define most avx instrs, with some exceptions.
% e.g. vadd:
%   | vadd lw nl, [a; b] => Some (interp_vector_binop Z.add (Z.of_N lw) 0 (N.to_nat nl) a b)



% The choice between these comes down to a tradeoff between rewriting rules and 
% If symbolic execution breaks an instr down into many simpler ops in the dag, then rewrite rules
% only need to normalize these existing ops (but must be powerful enough to see patterns).
% On the other hand, symbolic execution could emit a single atomic op like \verb|vadd| in the dag.
% This means that \verb|interp_op| needs to specify how this op decomposes in semantics, but also that
% rewrite rules need to handle this op in context.


% Then we get nice rewriting properties from lemmas like:
%
% (* Extracting lane k from a vector binop equals applying the scalar op to
%    the extracted lanes. Requires lo to be lane-aligned and in range. *)
% Lemma interp_vector_binop_slice_lane scalar_op lane_width lo sz a b num_lanes :
%   lane_width > 0 ->
%   sz = Z.to_N lane_width ->
%   (Z.of_N lo) mod lane_width = 0 ->
%   (lo + sz <= sz * num_lanes)%N ->
%   Z.land (scalar_op
%               (Z.land (Z.shiftr a (Z.of_N lo)) (Z.ones (Z.of_N sz)))
%               (Z.land (Z.shiftr b (Z.of_N lo)) (Z.ones (Z.of_N sz))))
%            (Z.ones (Z.of_N sz))
%   = Z.land (Z.shiftr (interp_vector_binop scalar_op lane_width 0 (N.to_nat num_lanes) a b)
%             (Z.of_N lo))
%          (Z.ones (Z.of_N sz)).
%
% ## Rewrite rules
%
% For example:
%
% (* Decompose a slice of a vector add into a scalar add of slices.
%    slice lo lw (vadd lw nl [v1; v2]) -> add lw [slice lo lw v1; slice lo lw v2]
%    when lo is lane-aligned (lo mod lw = 0) and in range (lo + lw <= lw * nl).
%    Uses coalescing_slice to collapse any nested slice(slice(...)) in the args,
%    since merge doesn't run rewrite passes on subexpressions. *)
% Definition slice_vadd (d : dag) :=
%   fun e => match e with
%     ExprApp (slice lo lw, [ExprApp (vadd lw' nl, [v1; v2])]) =>
%       if N.eqb lw lw' && N.eqb (lo mod lw)%N 0%N && N.leb (lo + lw) (lw' * nl) && N.ltb 0 lw
%       then ExprApp (add lw, [coalescing_slice lo lw v1; coalescing_slice lo lw v2])
%       else e | _ => e end%bool%N.
%
% Notice this requires a new helper rewrite function:
%
% (* Helper: build slice lo lw [v], coalescing nested slices.
%    If v = slice lo2 s2 [e'] and lo+lw <= s2, produce slice (lo2+lo) lw [e'] instead. *)
% Definition coalescing_slice (lo lw : N) (v : expr) : expr := ...
%
% since we get a lot of nested slices coming out of the slice_vadd rule otherwise,
% and the other rewrite rules aren't sophisticated enough to recurse into this.
%
% In fact, we ended up adding this rule as well:
%
% (* Recursively normalize (slice lo sz) over nested slice/set_slice towers,
%    peeling three kinds of layers at once:
%    - nested slice lo2 s2: combine offsets to (lo2+lo, sz, e')
%    - disjoint set_slice lo2 s2: descend into base
%    - containing set_slice lo2 s2: descend into val with adjusted lo
%    Needed for gather patterns (vmovq -> vpunpcklqdq -> vinserti128) that
%    interleave slice and set_slice layers, which the individual rules miss. *)
% Fixpoint slice_tower_normalize (lo sz : N) (inner : expr) (fuel : nat) : expr := ...
%
% Definition slice_tower (d : dag) :=
%   fun e => match e with
%     ExprApp (slice lo sz, [inner]) => slice_tower_normalize lo sz inner 16
%   | _ => e end.
%
% to handle similar patterns in the dag. Its unclear whether both are truly necessary,
% but one could check empirically.

